% Generated draft: requires mathematical and source review.
\documentclass[11pt]{article}
\usepackage[margin=25mm]{geometry}
\usepackage{fontspec}
\setmainfont{Latin Modern Roman}
\usepackage{amsmath,amssymb,mathtools,unicode-math}
\setmathfont{Latin Modern Math}
\usepackage{enumitem,needspace,xurl,hyperref}
\setlength{\parindent}{0pt}
\setlength{\parskip}{5pt}
\setlength{\emergencystretch}{3em}
\newcommand{\N}{\mathbb{N}}
\newcommand{\Z}{\mathbb{Z}}
\newcommand{\Q}{\mathbb{Q}}
\newcommand{\R}{\mathbb{R}}
\newcommand{\C}{\mathbb{C}}
\newcommand{\F}{\mathbb{F}}
\newcommand{\A}{\mathbb{A}}
\newcommand{\PP}{\mathbb{P}}
\newcommand{\E}{\mathbb{E}}
\newcommand{\eps}{\varepsilon}
\newcommand{\dd}{\,\mathrm{d}}
\begin{document}
\section*{UnsolvedMath ID: 30005736}
\textbf{Source problem code:} catalogue entry
\subsection{Definitions and mathematical statement}
The supplied record combines two requests whose relationship is not specified. Neither is sufficiently specified to yield a complete mathematical proposition from this record alone.

\textbf{Traffic control.} The intended model is a coupled traffic partial differential equation--ordinary differential equation system with an automated vehicle (AV) and a control. The excerpt refers to the Cauchy problem labelled (4)--(6), but does not supply those equations, their initial conditions, the road domain, boundary conditions, coupling conditions at the AV, or the admissible control class. It mentions a parameter \(\alpha\in(0,1)\), without specifying its role.

Here \(BV\) denotes bounded variation: on a specified open domain, an integrable function belongs to \(BV\) when its distributional derivative is a finite vector-valued Radon measure. The excerpt does not specify the domain, variables, local versus global interpretation, or precise solution class intended by its use of \(BV\). It reports an existence result for constant controls, for every initial condition in the same intended class, with solutions entropic outside the AV's location. The relevant entropy inequalities and the admissibility conditions at the AV are absent.

The requested extension is to establish existence for time-varying controls and to solve the stabilization problem for that same system. To express the existence target schematically, let \(\mathcal U_{\mathrm{tv}}\) denote the intended admissible time-varying controls, \(\mathcal B\) the intended initial-data class, and \(\mathcal S(u,z_0)\) the intended solutions of (4)--(6) with control \(u\) and initial datum \(z_0\), satisfying the required \(BV\) and entropy conditions. The requested assertion would have the form
\[
  \forall u\in\mathcal U_{\mathrm{tv}}\;\forall z_0\in\mathcal B,
  \qquad \mathcal S(u,z_0)\ne\varnothing.
\]
These symbols are placeholders for missing source definitions, not a specification of a new model. In particular, neither the admissible temporal regularity of \(u\) nor the time interval of existence is supplied. Stabilization additionally requires a specified target state or set, admissible controller information, a topology or error measure, and a convergence requirement; none is supplied. Thus a precise stabilization assertion cannot be recovered.

\textbf{Learning Lyapunov functions.} The excerpt describes a preliminary neural network trained to find Lyapunov functions for an unspecified nonlinear system and calls this a general open question. The requested task is to develop reliable neural methods for constructing such functions. The nonlinear dynamics, state space, equilibrium or target set, domain of certification, regularity and decrease requirements for the Lyapunov function, network class, and meaning of reliability are not supplied. Consequently, no particular existence, approximation, certification, or algorithmic guarantee can be stated without additional hypotheses.

The supplied literature assessment concerns a further formulation involving neural approximation of a drift for potentials described as strongly convex and Lipschitz. It reports a global linear error construction under \(M<\sqrt{2}\,m\), with \(m\), \(M\), the drift, and the error criterion undefined here. It explicitly says this does not establish the claim for all such potentials. The record does not explain how that formulation supplies partial progress on either requested target above.
\subsection{Short English statement}
Permanent UnsolvedMath ID \(30005736\) requests existence for time-varying controls and stabilization of an unspecified traffic PDE--ODE system, together with reliable neural construction of Lyapunov functions. The referenced equations and essential hypotheses are missing, so these remain qualified research requests rather than complete mathematical propositions. The supplied partially solved label and literature assessment require independent review.
\subsection{Sources}
\begin{itemize}
\item Supplied local UnsolvedMath registry record: permanent ID \(30005736\), title \emph{Time-Varying Traffic Control and Lyapunov Learning}, exact source problem code \texttt{catalogue entry}, in the Oberwolfach Reports Open Problems collection. The record lists proposed year 2023 and supplies no proposer attribution.
\item Supplied report citation: \emph{Control Methods in Hyperbolic PDEs} (2024), Oberwolfach Reports, pp. 2941--3011, DOI 10.4171/owr/2023/52. \url{https://doi.org/10.4171/owr/2023/52}. The report is described as documenting mathematical workshops and open-problem discussions. The supplied literature list also labels this source \emph{Oberwolfach Report 52/2023}; the 2024 citation year and 2023 report designation are retained as supplied.
\item Supplied literature source: Eigel, Miranda, Schuette, and Sommer, WIAS Preprint 3077. \url{https://www.wias-berlin.de/preprint/3077/wias_preprints_3077.pdf}. The supplied assessment reports neural networks with global linear error for a drift under the additional restriction \(M<\sqrt{2}\,m\), and explicitly limits this to partial progress rather than a proof for all strongly convex Lipschitz potentials. The relevant definitions and its connection to this record's targets are missing.
\item Literature assessment date: \texttt{2026-08-22}, copied exactly from \texttt{literature\_checked\_at}. The supplied status is \texttt{partially\_solved}. This is supplied dated curation derived from the OWR material, not an independently verified literature assessment; its status and cited claims require independent verification.
\item The original excerpt attributes constant-control existence to reference [5] and field experiments to reference [4], without supplying either bibliographic entry. It reports experiments in Nashville, Tennessee, in November 2022, within the CIRCLES project, involving 100 autonomous cars on a highway in dense peak-hour traffic, and reports that the candidate controller drastically reduced the speed variance of stop-and-go waves. These are attributed source reports, not verified mathematical stabilization guarantees.
\item The registry labels the statement \texttt{corrected\_verified} and says it was checked against the official report and preserves its hypotheses, notation, and target. Those are supplied verification claims; no source was independently visited or verified in preparing this draft, and the supplied excerpt omits essential mathematical data.
\end{itemize}
\end{document}
